% GENERATED FILE. Edit canonical lessons, then run npm run generate:books.
% canonical-source-hash: fnv1a64:1726ff595cad512f
% canonical-lessons: ES-C524-incendio, ES-C524-incendiar, ES-C524-arder, ES-C524-alarma, ES-C524-evacuar

\chapter{El incendio}
\label{ch:es-el-incendio}
\hlchaptermodality{524}

\begin{chapteropening}
\textbf{By the end of this chapter:} \emph{I can report a fire: what burned, the alarm, and the people who were evacuated.}

Complete the last lesson of chapter 524: I can report a fire: what burned, the alarm, and the people who were evacuated.
\end{chapteropening}

\section[el incendio]{el incendio — the fire (big, destructive)}

\label{lesson:ES-C524-incendio}



\begin{quote}
Before the new one: say the Spanish for to testify, then the Spanish for to free.
\end{quote}

\subsection*{You'll want to know: el incendio}
\textbf{el incendio} — \textquotedblleft{}the fire (big, destructive)\textquotedblright{}.

\begin{quote}
El incendio empezó en la cocina.
\end{quote}

\begin{quote}
\emph{The fire started in the kitchen.}
\end{quote}

\emph{El fuego} is fire in general; \emph{un incendio} is a fire that burns a building or a forest.

\begin{cousinweb}
\emph{incendio} is a learned borrowing from Latin \emph{incendium} 'a fire, a blaze' (EtymDB), built on \emph{incendere} 'to set on fire'. English \emph{incendiary} comes from the same verb.
\end{cousinweb}

\begin{grammarlens}[title={What you've built}]
One more word for telling, explaining and describing at B1.
\end{grammarlens}

\subsection*{Your turn}
\begin{itemize}
\raggedright
  \item \emph{Say it:} \emph{el incendio}
  \item \emph{Say it:} \emph{el incendio}, once more
  \item \emph{Say it:} a sentence of your own with \emph{el incendio}
  \item \emph{Recall:} say the Spanish for to testify, then the Spanish for to free, then say \emph{el incendio} again
\end{itemize}

\begin{tcolorbox}[breakable,colback=teal!4,colframe=teal!35!black,title={Before you move on}]
What does el incendio mean? (\textquotedblleft{}The fire (big, destructive)\textquotedblright{}.) Say it once more.
\end{tcolorbox}
\section[incendiar]{incendiar — to set on fire}

\label{lesson:ES-C524-incendiar}



\begin{quote}
Before the new one: say the Spanish for to free, then the Spanish for the fire.
\end{quote}

\subsection*{You'll want to know: incendiar}
\textbf{incendiar} — \textquotedblleft{}to set on fire\textquotedblright{}.

\begin{quote}
Alguien incendió el coche de noche.
\end{quote}

\begin{quote}
\emph{Someone set the car on fire at night.}
\end{quote}

\emph{Incendiarse} is to catch fire: \emph{La casa se incendió.}

\begin{cousinweb}
\emph{incendiar} is built on \emph{incendio} 'a fire', a learned word from Latin \emph{incendium}, built on \emph{incendere} 'to set on fire' (Lewis \& Short).
\end{cousinweb}

\begin{grammarlens}[title={What you've built}]
One more word for telling, explaining and describing at B1.
\end{grammarlens}

\subsection*{Your turn}
\begin{itemize}
\raggedright
  \item \emph{Say it:} \emph{incendiar}
  \item \emph{Say it:} \emph{incendiar}, once more
  \item \emph{Say it:} a sentence of your own with \emph{incendiar}
  \item \emph{Recall:} say the Spanish for to free, then the Spanish for the fire, then say \emph{incendiar} again
\end{itemize}

\begin{tcolorbox}[breakable,colback=teal!4,colframe=teal!35!black,title={Before you move on}]
What does incendiar mean? (\textquotedblleft{}To set on fire\textquotedblright{}.) Say it once more.
\end{tcolorbox}
\section[arder]{arder — to burn, to be on fire}

\label{lesson:ES-C524-arder}



\begin{quote}
Before the new one: say the Spanish for the fire, then the Spanish for to set on fire.
\end{quote}

\subsection*{You'll want to know: arder}
\textbf{arder} — \textquotedblleft{}to burn, to be on fire\textquotedblright{}.

\begin{quote}
El edificio ardió toda la noche.
\end{quote}

\begin{quote}
\emph{The building burned all night.}
\end{quote}

\emph{Arder} is to be burning; to burn something is \emph{quemar}.

\begin{cousinweb}
\emph{arder} is inherited from Latin \emph{ārdēre} 'to be on fire, to blaze' (EtymDB). English \emph{ardent} comes from the same verb.
\end{cousinweb}

\begin{grammarlens}[title={What you've built}]
One more word for telling, explaining and describing at B1.
\end{grammarlens}

\subsection*{Your turn}
\begin{itemize}
\raggedright
  \item \emph{Say it:} \emph{arder}
  \item \emph{Say it:} \emph{arder}, once more
  \item \emph{Say it:} a sentence of your own with \emph{arder}
  \item \emph{Recall:} say the Spanish for the fire, then the Spanish for to set on fire, then say \emph{arder} again
\end{itemize}

\begin{tcolorbox}[breakable,colback=teal!4,colframe=teal!35!black,title={Before you move on}]
What does arder mean? (\textquotedblleft{}To burn, to be on fire\textquotedblright{}.) Say it once more.
\end{tcolorbox}
\section[la alarma]{la alarma — the alarm}

\label{lesson:ES-C524-alarma}



\begin{quote}
Before the new one: say the Spanish for to set on fire, then the Spanish for to burn.
\end{quote}

\subsection*{You'll want to know: la alarma}
\textbf{la alarma} — \textquotedblleft{}the alarm\textquotedblright{}.

\begin{quote}
Sonó la alarma a las tres de la mañana.
\end{quote}

\begin{quote}
\emph{The alarm went off at three in the morning.}
\end{quote}

\emph{Dar la alarma} — to raise the alarm.

\begin{cousinweb}
\emph{alarma} comes from the Italian cry \emph{all'arme} 'to arms' (EtymDB), and \emph{arme} goes back to Latin \emph{arma} 'weapons'. An alarm was first a call to take up your weapons.
\end{cousinweb}

\begin{grammarlens}[title={What you've built}]
One more word for telling, explaining and describing at B1.
\end{grammarlens}

\subsection*{Your turn}
\begin{itemize}
\raggedright
  \item \emph{Say it:} \emph{la alarma}
  \item \emph{Say it:} \emph{la alarma}, once more
  \item \emph{Say it:} a sentence of your own with \emph{la alarma}
  \item \emph{Recall:} say the Spanish for to set on fire, then the Spanish for to burn, then say \emph{la alarma} again
\end{itemize}

\begin{tcolorbox}[breakable,colback=teal!4,colframe=teal!35!black,title={Before you move on}]
What does la alarma mean? (\textquotedblleft{}The alarm\textquotedblright{}.) Say it once more.
\end{tcolorbox}
\section[evacuar]{evacuar — to evacuate}

\label{lesson:ES-C524-evacuar}



\begin{quote}
Before the new one: say the Spanish for to burn, then the Spanish for the alarm.
\end{quote}

\subsection*{You'll want to know: evacuar}
\textbf{evacuar} — \textquotedblleft{}to evacuate\textquotedblright{}.

\begin{quote}
Evacuaron el edificio en diez minutos.
\end{quote}

\begin{quote}
\emph{They evacuated the building in ten minutes.}
\end{quote}

\begin{cousinweb}
\emph{evacuar} is a learned borrowing from Latin \emph{ēvacuāre} 'to empty out' (EtymDB), built on \emph{vacuus} 'empty'. To evacuate a building is to empty it.
\end{cousinweb}

\begin{grammarlens}[title={What you've built}]
One more word for telling, explaining and describing at B1.
\end{grammarlens}

\subsection*{Your turn}
\begin{itemize}
\raggedright
  \item \emph{Say it:} \emph{evacuar}
  \item \emph{Say it:} \emph{evacuar}, once more
  \item \emph{Say it:} a sentence of your own with \emph{evacuar}
  \item \emph{Recall:} say the Spanish for to burn, then the Spanish for the alarm, then say \emph{evacuar} again
\end{itemize}

\begin{tcolorbox}[breakable,colback=teal!4,colframe=teal!35!black,title={Before you move on}]
What does evacuar mean? (\textquotedblleft{}To evacuate\textquotedblright{}.) Say it once more.
\end{tcolorbox}
